% problem_id: S001-spaces-chow-lemma-key-formula
% source_id: S001 (Stacks Project)
% source_file: spaces-chow.tex
% statement_locator: spaces-chow.tex lines 2867--2880
% proof_locator: spaces-chow.tex lines 2882--3063
% env: lemma
% id_note: label 跨文件重名 ⇒ id 用文件限定形式
% pairing: tight (verified)
% status: complete (statement + author proof verbatim + reason + locators)
%
\begin{lemma}[Key formula]
\label{lemma-key-formula}
In the situation above the cycle
$$
\sum
(Z_i \to X)_*\left(
\text{ord}_{B_i}(f_i) \text{div}_{\mathcal{N}|_{Z_i}}(t_i|_{Z_i}) -
\text{ord}_{B_i}(g_i) \text{div}_{\mathcal{L}|_{Z_i}}(s_i|_{Z_i}) \right)
$$
is equal to the cycle
$$
\sum (Z_i \to X)_*\text{div}(\partial_{B_i}(f_i, g_i))
$$
\end{lemma}

\begin{proof}
The strategy of the proof will be: first reduce to the case where
$\mathcal{L}$ and $\mathcal{N}$ are trivial invertible modules,
then change our choices of local trivializations, and then finally
use \'etale localization to reduce to the case of schemes\footnote{It
is possible that a shorter proof can be given by immediately
applying \'etale localization.}.

\medskip\noindent
First step. Let $q : T \to X$ be the morphism constructed in
Lemma \ref{lemma-Gm-torsor}. We will use all properties stated in
that lemma without further mention. In particular, it suffices to
show that the cycles are equal after pulling back by $q$.
Denote $s'$ and $t'$ the pullbacks of $s$ and $t$ to meromorphic sections
of $q^*\mathcal{L}$ and $q^*\mathcal{N}$.
Denote $Z'_i = q^{-1}(Z_i)$, denote $\xi'_i \in |Z'_i|$ the generic point,
denote $B'_i = \mathcal{O}_{T, \xi'_i}^h$, denote
$\mathcal{L}_{\xi'_i}$ and $\mathcal{N}_{\xi'_i}$ the pullbacks
of $\mathcal{L}$ and $\mathcal{N}$ to $\Spec(B'_i)$.
Recall that we have commutative diagrams
$$
\xymatrix{
\Spec(B'_i) \ar[r]_-{c_{\xi'_i}} \ar[d] & T \ar[d]^q \\
\Spec(B_i) \ar[r]^-{c_{\xi_i}} & X
}
$$
see Decent Spaces, Remark
\ref{decent-spaces-remark-functoriality-henselian-local-ring}.
Denote $s'_i$ and $t'_i$ the pullbacks of $s_i$ and $t_i$ which
are generators of $\mathcal{L}_{\xi'_i}$ and $\mathcal{N}_{\xi'_i}$.
Then we have
$$
s' = f'_i s'_i
\quad\text{and}\quad
t' = g'_i t'_i
$$
where $f'_i$ and $g'_i$ are the images of $f_i, g_i$
under the map $Q(B_i) \to Q(B'_i)$ induced by $B_i \to B'_i$.
By Algebra, Lemma \ref{algebra-lemma-pullback-module} we have
$$
\text{ord}_{B_i}(f_i) = \text{ord}_{B'_i}(f'_i)
\quad\text{and}\quad
\text{ord}_{B_i}(g_i) = \text{ord}_{B'_i}(g'_i)
$$
By Lemma \ref{lemma-prepare-flat-pullback-cap-c1} applied
to $q : Z'_i \to Z_i$ we have
$$
q^*\text{div}_{\mathcal{N}|_{Z_i}}(t_i|_{Z_i}) =
\text{div}_{q^*\mathcal{N}|_{Z'_i}}(t'_i|_{Z'_i})
\quad\text{and}\quad
q^*\text{div}_{\mathcal{L}|_{Z_i}}(s_i|_{Z_i}) =
\text{div}_{q^*\mathcal{L}|_{Z'_i}}(s'_i|_{Z'_i})
$$
This already shows that the first cycle in the statement of the
lemma pulls back to the corresponding cycle for
$s', t', Z'_i, s'_i, t'_i$. To see the same is true for the
second, note that by
Chow Homology, Lemma \ref{chow-lemma-tame-symbol-formally-smooth} we have
$$
\partial_{B_i}(f_i, g_i) \mapsto
\partial_{B'_i}(f'_i, g'_i)
\quad\text{via}\quad
\kappa(\xi_i) \to \kappa(\xi'_i)
$$
Hence the same lemma as before shows that
$$
q^*\text{div}(\partial_{B_i}(f_i, g_i)) =
\text{div}(\partial_{B'_i}(f'_i, g'_i))
$$
Since $q^*\mathcal{L} \cong \mathcal{O}_T$ we find that it suffices
to prove the equality in case $\mathcal{L}$ is trivial.
Exchanging the roles of $\mathcal{L}$ and $\mathcal{N}$ we
see that we may similarly assume $\mathcal{N}$ is trivial.
This finishes the proof of the first step.

\medskip\noindent
Second step. Assume $\mathcal{L} = \mathcal{O}_X$ and
$\mathcal{N} = \mathcal{O}_X$. Denote $1$ the trivializing
section of $\mathcal{L}$. Then $s_i = u \cdot 1$
for some unit $u \in B_i$. Let us examine what happens
if we replace $s_i$ by $1$. Then $f_i$ gets replaced by $u f_i$.
Thus the first part of the first expression of the lemma is unchanged
and in the second part we add
$$
\text{ord}_{B_i}(g_i)\text{div}(u|_{Z_i})
$$
where $u|_{Z_i}$ is the image of $u$ in the residue field by
Spaces over Fields, Lemma
\ref{spaces-over-fields-lemma-divisor-meromorphic-well-defined}
and in the second expression we add
$$
\text{div}(\partial_{B_i}(u, g_i))
$$
by bi-linearity of the tame symbol. These terms agree by the property
of the tame symbol given in
Chow Homology, Equation (\ref{chow-item-normalization}).

\medskip\noindent
Let $Y \subset X$ be an integral closed subspace with $\dim_\delta(Y) = n - 2$.
To show that the coefficients of $Y$ of the two cycles of the lemma
is the same, we may do a replacement of $s_i$ by $1$ as in the previous
paragraph. In exactly the same way one shows that we may do a replacement
of $t_i$ by $1$. Since there are only a finite number of $Z_i$
such that $Y \subset Z_i$ we may assume $s_i = 1$ and $t_i = 1$
for all these $Z_i$.

\medskip\noindent
Third step. Here we prove the coefficients of $Y$ in the cycles
of the lemma agree for an integral closed subspace $Y$ with
$\dim_\delta(Y) = n - 2$ such that moreover
$\mathcal{L} = \mathcal{O}_X$ and $\mathcal{N} = \mathcal{O}_X$
and $s_i = 1$ and $t_i = 1$ for all $Z_i$ such that $Y \subset Z_i$.
After replacing $X$ by a smaller open subspace we may
in fact assume that $s_i$ and $t_i$ are equal to $1$ for all $i$.
In this case the first cycle is zero. Our task is to show that
the coefficient of $Y$ in the second cycle is zero as well.

\medskip\noindent
First, since $\mathcal{L} = \mathcal{O}_X$ and $\mathcal{N} = \mathcal{O}_X$
we may and do think of $s, t$ as rational functions $f, g$ on $X$.
Since $s_i$ and $t_i$ are equal to $1$ we find that $f_i$, resp.\ $g_i$
is the image of $f$, resp.\ $g$ in $Q(B_i)$ for all $i$.
Let $\zeta \in |Y|$ be the generic point. Choose an
\'etale neighbourhood
$$
(U, u) \longrightarrow (X, \zeta)
$$
and denote $Y' = \overline{\{u\}} \subset U$.
Since an \'etale morphism is flat, we can pullback $f$ and $g$
to regular meromorphic functions on $U$ which we will also denote
$f$ and $g$.
For every prime divisor $Y \subset Z \subset X$ the scheme $Z \times_X U$
is a union of prime divisors of $U$. Conversely, given a prime divisor
$Y' \subset Z' \subset U$, there is a prime divisor
$Y \subset Z \subset X$ such that $Z'$ is a component of $Z \times_X U$.
Given such a pair $(Z, Z')$ the ring map
$$
\mathcal{O}_{X, \xi}^h \to \mathcal{O}_{U, \xi'}^h
$$
is \'etale (in fact it is finite \'etale). Hence we find that
$$
\partial_{\mathcal{O}_{X, \xi}^h}(f, g) \mapsto
\partial_{\mathcal{O}_{U, \xi'}^h}(f, g)
\quad\text{via}\quad
\kappa(\xi) \to \kappa(\xi')
$$
by Chow Homology, Lemma \ref{chow-lemma-tame-symbol-formally-smooth}.
Thus Lemma \ref{lemma-etale-pullback-principal-divisor} applies to show
$$
(Z \times_X U \to Z)^*\text{div}_Z(\partial_{\mathcal{O}_{X, \xi}^h}(f, g))
=
\sum\nolimits_{Z' \subset Z \times_X U}
\text{div}_{Z'}(\partial_{\mathcal{O}_{U, \xi'}^h}(f, g))
$$
Since flat pullback commutes with pushforward along closed
immersions (Lemma \ref{lemma-flat-pullback-proper-pushforward})
we see that it suffices to prove that the coefficient
of $Y'$ in
$$
\sum\nolimits_{Z' \subset U} 
(Z' \to U)_*\text{div}_{Z'}(\partial_{\mathcal{O}_{U, \xi'}^h}(f, g))
$$
is zero.

\medskip\noindent
Let $A = \mathcal{O}_{U, u}$. Then $f, g \in Q(A)^*$.
Thus we can write $f = a/b$ and $g = c/d$ with
$a, b, c, d \in A$ nonzerodivisors.
The coefficient of $Y'$ in the expression above is
$$
\sum\nolimits_{\mathfrak q \subset A\text{ height }1}
\text{ord}_{A/\mathfrak q}(\partial_{A_\mathfrak q}(f, g))
$$
By bilinearity of $\partial_A$ it suffices to prove
$$
\sum\nolimits_{\mathfrak q \subset A\text{ height }1}
\text{ord}_{A/\mathfrak q}(\partial_{A_\mathfrak q}(a, c))
$$
is zero and similarly for the other pairs $(a, d)$, $(b, c)$, and
$(b, d)$. This is true by
Chow Homology, Lemma \ref{chow-lemma-key-nonzerodivisors}.
\end{proof}
